%!TEX root = ../LLMSeminar.tex
% ──────────────────────────────────────────────────────────────
%  Section 12 — Appendix: Annotated reference coding agent
% ──────────────────────────────────────────────────────────────

\appendix
\section{Annotated reference coding agent}\label{sec:appendix-agent}

For completeness, we record here a pedagogically transparent
implementation of the loop discussed in
\S\ref{sec:agent-loop}.  The code is intended to be read,
not optimised; production agents add many features (stream
responses, retry on tool failure, manage context with
summarisation, sandbox execution, log to disk) but none
that change the core architecture.

This implementation uses the Anthropic Python SDK and the
\texttt{claude-sonnet-4-6} model; the structure is identical
for OpenAI's or Google's APIs with minor adjustments to the
schema names.

\subsection{Tools schema}

\begin{lstlisting}
TOOLS = [
  {
    "name": "read_file",
    "description": "Read a UTF-8 text file and return its contents.",
    "input_schema": {
      "type": "object",
      "properties": { "path": { "type": "string" } },
      "required": ["path"],
    },
  },
  {
    "name": "write_file",
    "description": "Overwrite a UTF-8 text file with the given content.",
    "input_schema": {
      "type": "object",
      "properties": {
        "path":    { "type": "string" },
        "content": { "type": "string" },
      },
      "required": ["path", "content"],
    },
  },
]
\end{lstlisting}

Each tool is described by a name (the LLM will reference it
by exactly this string), a one-sentence \texttt{description}
(the model uses this to decide \emph{when} to call the tool),
and a JSON-Schema \texttt{input\_schema} fixing the shape
of the arguments.

\subsection{Tool dispatcher}

\begin{lstlisting}
def run_tool(name, args):
    if name == "read_file":
        with open(args["path"], "r") as f:
            return f.read()
    if name == "write_file":
        os.makedirs(os.path.dirname(args["path"]) or ".", exist_ok=True)
        with open(args["path"], "w") as f:
            f.write(args["content"])
        return f"wrote {len(args['content'])} bytes to {args['path']}"
    raise ValueError(f"unknown tool: {name}")
\end{lstlisting}

The dispatcher is plain Python.  No AI involvement; it does
exactly what the agent asked, no more.

\subsection{System prompt}

\begin{lstlisting}
SYSTEM = """
You are a coding agent with read- and write-file tools.
Your own source code lives at agent.py.  You may read and
write it to improve yourself.  Changes take effect on next
invocation.  Preserve the loop and tool plumbing in this
script.
"""
\end{lstlisting}

The system prompt establishes the agent's identity and the
boundaries within which it should self-modify.  It is short
on purpose.  Long system prompts in pedagogical examples
encourage the wrong instinct---``promptcraft as central
art''---when in fact the loop is what does the work.

\subsection{The loop}

\begin{lstlisting}
def run_agent(user_input, max_turns=25):
    messages = [{"role": "user", "content": user_input}]

    for _ in range(max_turns):
        response = client.messages.create(
            model       = "claude-sonnet-4-6",
            max_tokens  = 4096,
            system      = SYSTEM,
            tools       = TOOLS,
            messages    = messages,
        )

        # 1. Append the assistant's reply to history.
        messages.append({"role": "assistant", "content": response.content})

        # 2. Stop if the model did not request any tool calls.
        if response.stop_reason != "tool_use":
            break

        # 3. Otherwise, run each requested tool and append results.
        tool_results = []
        for block in response.content:
            if block.type == "tool_use":
                try:
                    result = run_tool(block.name, block.input)
                except Exception as e:
                    result = f"ERROR: {e}"
                tool_results.append({
                    "type": "tool_result",
                    "tool_use_id": block.id,
                    "content": result,
                })

        messages.append({"role": "user", "content": tool_results})

    return messages

if __name__ == "__main__":
    run_agent("Read your own source code and improve yourself.")
\end{lstlisting}

\subsection{Reading the listing}

A few features are worth dwelling on.

\paragraph{The \texttt{role} key.}
Every message has a \texttt{role}---\texttt{user},
\texttt{assistant}, or, when bundling tool outputs, \texttt{user}
again with structured \texttt{tool\_result} blocks inside.
The role is what allows the LLM to predict the next token in
context.  Without it, the model would have to guess; with it,
training-time conditioning takes over and the model behaves
as expected.

\paragraph{The \texttt{tool\_use\_id}.}
When the LLM emits a tool call it is given a unique
\texttt{tool\_use\_id}; the corresponding result must reference
the same id.  This pairing lets the LLM line up requests and
answers when there are several tool calls in flight.

\paragraph{The \texttt{stop\_reason}.}
The vendor returns a discriminator indicating why the model
stopped generating.  Values include \texttt{end\_turn} (the
model is finished), \texttt{tool\_use} (it has emitted a
tool-call), \texttt{max\_tokens} (it ran out of budget).  The
loop dispatches on this.

\paragraph{The cap \texttt{max\_turns}.}
A safety latch.  Without one, an agent given a vague task
will gladly recurse for hours and consume significant API
credit.  Twenty-five turns is enough for almost any
self-bootstrapping experiment; raise it deliberately.

\subsection{Self-improvement, observed}

Running this agent with the prompt
\emph{``read your own source code and improve yourself''}
produces a recognisable trajectory.  In one run observed in
preparation:

\begin{enumerate}[(1)]
  \item The agent reads \texttt{agent.py}.
  \item It writes a comment-rich version of itself, adds a
    docstring, and refactors \texttt{run\_tool} into a dict
    dispatch.
  \item It declares an \texttt{exec\_bash} tool, including
    its JSON schema.
  \item It writes back \texttt{agent.py}.
  \item On the next invocation, it has bash; it then asks
    what to do, and proposes useful self-experiments.
\end{enumerate}

Total wall clock: under ten minutes.  Total API cost: a few
euros.

\begin{remark}
  This experiment is genuinely worth running.  It is the
  shortest path from understanding the loop on paper to
  understanding what it feels like to be a participant in
  the loop in the wild.  The \emph{strangeness} of seeing an
  agent declare a new capability for itself, write the code,
  and then exercise it---this is the lived form of the
  technical content of \S\ref{sec:agent-loop}--%
  \S\ref{sec:self-improvement}.
\end{remark}

\subsection{Cost and safety}

\begin{itemize}
  \item Run in a Docker container or a fresh user account.
    A read/write/exec agent given network access can do
    surprising things, and ``surprising'' includes
    ``destructive.''
  \item Set a hard quota on API spend.  All vendors offer
    this in their billing dashboard.
  \item Log the full message stream to disk.  Reading a
    transcript afterwards is the most efficient way to learn
    what your agent is actually doing.
  \item Resist the urge to over-engineer the system prompt.
    The loop, the tools, and feedback are what produce
    capability; the prompt mostly steers.
\end{itemize}
